\chapter{Implementation: SCION Path-Selection Pipeline}
\label{ch:p1impl}

This chapter describes how the Part~I design is realized: the agent family and their architectures
(\cref{sec:p1impl:agents}), the end-to-end evaluation pipeline that turns a topology into trained models
and comparable results (\cref{sec:p1impl:pipeline}), and the baselines and probing accounting against
which the learned selectors are measured (\cref{sec:p1impl:baselines}). The implementation is in Python
(PyTorch for the agents) and is organized so that each pipeline stage passes data to the next through
timestamped run directories.

\section{Agent Family}
\label{sec:p1impl:agents}

The codebase implements a family of agents of increasing capability, which also serve as the ablation
ladder in \cref{ch:p1eval}.

\begin{description}
  \item[Flat DQN (simple / enhanced).] A conventional fixed-action DQN over the 5-dimensional global
  state. The \emph{simple} variant is a small multilayer perceptron with $\varepsilon$-greedy
  exploration and a periodic target network; the \emph{enhanced} variant adds Double
  DQN~\cite{VanHasselt2016}, a dueling head~\cite{Wang2016}, prioritized experience
  replay~\cite{Schaul2016}, soft target updates, and action masking. These agents cannot handle a
  variable path set and stand in for the original design~\cite{Bourak2025}.
  \item[Path-scoring DQN (simple / enhanced).] The permutation-equivariant design of
  \cref{sec:p1design:scoring}: a shared per-path scorer producing one $Q$-value per path, with variable
  $N$ handled by padding to the batch maximum and masking. The enhanced variant uses the dueling scoring
  head, Double DQN, and prioritized replay. This family delivers the dynamic-action property.
  \item[Conditional path-scoring DQN.] The recommended agent, adding the weight-FiLM dueling head of
  \cref{sec:p1design:film}. Its global state is extended with the five-dimensional encoded intent vector,
  and the FiLM sub-network modulates the per-path features feeding the advantage stream. Two intent
  encodings are supported and recorded in the checkpoint -- the raw weights, or a policy-friendly rescaling
  to $[-1,1]$ for smoother gradients.
\end{description}

\paragraph{Training details.} Prioritized replay uses priority exponent $\alpha = 0.6$ with the
importance-sampling exponent $\beta$ annealed from $0.4$ to $1.0$. Target networks use soft (Polyak)
updates. Invalid actions are masked to a large negative value both when acting and when forming
bootstrap targets, so masked paths never contribute to the loss. The conditional agent is trained across
reward profiles on a stratified schedule (\cref{sec:p1design:film}), so that within a single training run
the network sees every intent an equal number of times.

\section{Evaluation Pipeline}
\label{sec:p1impl:pipeline}

The system is driven by a six-stage pipeline, each stage writing artifacts consumed by the next:

\begin{enumerate}
  \item \textbf{Topology generation} -- BRITE produces an AS graph, converted into SCION roles, ISDs, and
  links.
  \item \textbf{Beaconing} -- PCB propagation discovers up/core/down segments and composes candidate paths
  per source--destination pair.
  \item \textbf{Traffic simulation} -- 28~days of hourly demand with per-link aggregation, producing the
  time-varying link states the environment replays.
  \item \textbf{Training} -- all agent variants are trained on the resulting environment; checkpoints and
  per-episode statistics are saved.
  \item \textbf{Evaluation} -- every method (learned and baseline) is run over held-out pairs, hours, and
  intents, logging reward, achieved latency and bandwidth, probing overhead, and selection time.
  \item \textbf{Figures} -- results are rendered in a consistent style for the report.
\end{enumerate}

Running the whole pipeline is a single command; intermediate artifacts (topology, path store, link
states, checkpoints, results tables) live under a timestamped run directory so that any stage can be
re-run or inspected in isolation. Key knobs -- topology size, number of training episodes, the cap on
training pairs, and learning hyperparameters -- are exposed as configuration so that experiments can be
scaled up or down.

\section{Baselines and Probing Accounting}
\label{sec:p1impl:baselines}

\paragraph{Baselines.} Six heuristic selectors provide the comparison points: shortest-path,
widest-path, lowest-latency, ECMP, the SCION default selection, and random. They span the spectrum from
naive (random) to strong single-metric heuristics (widest for throughput, lowest-latency for delay), so
that the learned selectors must beat not just a weak baseline but the \emph{right} heuristic for each
intent.

\paragraph{Probing accounting.} A fair comparison must charge each method for the measurements it needs.
The environment therefore records, per decision, which paths a method probed and at what cost -- latency
probes being cheap and full bandwidth probes expensive (\cref{sec:p1design:env}). Heuristics that require
a metric on every path (e.g.\ widest-path) incur the cost of probing all paths, whereas the learned
selector is charged only for the paths it actually inspects. This makes the probing-overhead reduction a
directly measurable quantity rather than an assumed benefit, and is the mechanism by which we reproduce
and test the near-oracle-quality-at-low-probing claim of Bourak~et~al.~\cite{Bourak2025}.
